\documentclass[11pt]{article}
\usepackage[margin=.86in]{geometry}
\usepackage{amsmath,amssymb,amsthm,mathtools}
\usepackage[hidelinks]{hyperref}
\newtheorem{theorem}{Theorem}
\newtheorem{lemma}{Lemma}
\newcommand{\E}{\mathbb E}
\newcommand{\PP}{\mathbb P}
\newcommand{\1}{\mathbf{1}}
\newcommand{\ph}{\varphi}
\newcommand{\Sg}{\mathcal S}
\title{Poisson Limits for Primes in Moving Arithmetic Progressions\\
\large Density-One Integer Moduli and Prescribed Patterns at Prime Moduli}
\author{Keio\\\small Combined manuscript from prime tournament rounds 9 and 10}
\date{November 2026}
\begin{document}
\maketitle
\begin{abstract}
We prove two unconditional Poisson limits for primes in moving
arithmetic progressions. On a natural-density-one set of integer moduli
$q$, the prime locations in a uniformly chosen reduced class, scaled
by $\varphi(q)\log q$, approach a unit-rate Poisson point process.
In particular the first-prime waiting time has an exponential limit.
On a relative-density-one set of prime moduli, one may instead choose
uniformly among the classes for which all rows of a prescribed
fixed-size admissible pattern are prime. The fraction of such classes
has an asymptotic involving the pattern's singular series, and after
the prescribed primes are removed the remaining process is still
Poisson. The pattern may move between dyadic blocks, but is chosen
before the modulus is sampled. Both proofs use a second residue
representative at the same modulus to establish concentration.
\end{abstract}

\section{Setting and relation between the results}
At a fixed modulus $q$ the randomness is in the choice of the
representative $a$, not in the primes. The points in the process are
rescaled prime \emph{values}; waiting times are the first point's
location and subsequent gaps between ordered points. In
Theorem~\ref{thm:integer}, the good moduli have natural density one
among all positive integers. In Theorem~\ref{thm:prime}, they have
relative density one among the primes, and the distribution of $a$ is
restricted to classes satisfying a prime-pattern event. Neither
result is a theorem for every modulus. In particular,
Theorem~\ref{thm:integer} need not apply to prime moduli, and ordinary
weak convergence does not justify conditioning on a sequence of rare
events.

The shared analytic input is Bienvenu's theorem for affine forms with
polylogarithmic coefficients \cite{bien}. For first and second moments,
one uses one and two independent representatives, respectively, while
keeping the same modulus. In the integer-modulus argument the local
factors distinguish $q=0$ from $q\ne0$ modulo each small prime. In the
prime-modulus argument, adding $q$ itself as a prime form removes the
zero residue at each local factor and permits direct normalization by
the prescribed pattern's singular series. The two full arguments
follow separately.

\section{The first-prime process at integer moduli}

For each $q$ let $a$ be uniform in $(\mathbb Z/q\mathbb Z)^\times$. Write
\[
 \Xi_q=\sum_{p\ {\rm prime},\ p\equiv a\ (q)}\delta_{p/(\ph(q)\log q)}.
\]
Point-process convergence is distributional for the vague topology on $[0,\infty)$; the only random variable at fixed $q$ is $a$.
\begin{theorem}\label{thm:integer}
There exists a single set $\mathcal Q$ of natural density one such that, as $q\to\infty$ in $\mathcal Q$, $\Xi_q$ converges to the unit-rate Poisson point process. In particular on any fixed disjoint bounded intervals $I_1,\ldots,I_t$ the limiting counts are independent $\operatorname{Pois}(|I_j|)$, and the first $k$ waiting times from zero for any fixed $k$ have joint limit $\operatorname{Exp}(1)^{\otimes k}$.
\end{theorem}
Leung \cite{lecounts,lepp} gives the full growing-modulus counterpart assuming a uniform prime-tuples hypothesis. The quantifier in Theorem~\ref{thm:integer} is weaker, but no unproved hypothesis enters it. Natural density one does not include every prescribed thin sequence of moduli, in particular not necessarily prime or primorial moduli. This uses the classical averaged-modulus idea of allowing the prime forms to have an additional independent variable. The essential available analytic ingredient is Bienvenu's polylogarithmic-coefficient linear-forms theorem \cite{bien}; bounded-coefficient Green--Tao--Ziegler alone would not suffice. Example 1.1 in \cite{bien} already counts sparse rows of arbitrary progressions, without a fixed-modulus empirical-process conclusion. We make no claim about the extreme largest least prime.

\subsection{Raw row moments with a common modulus}
Throughout one dyadic block $N<q\le2N$ write $L=\log N$, $u(q)=q/\ph(q)$. Fix $r\ge1$ and a box $\mathcal I=I_1\times\cdots\times I_r$ of bounded intervals of positive lengths in $[0,\infty)$; they need not be disjoint. Let $J=\prod |I_i|$. For each $i$ range $b_i$ over positive integers in $L I_i$ and write a star for pairwise distinct $b_i$. Define
\[
 Z_q=\frac1q\sum_{1\le a\le q}\sum_{\mathbf b}^{*}
     \prod_{i=1}^r\1_{\mathbb P}(a+b_i q).
\]
This is an average over \emph{all} representatives. If any product is one then $(a,q)=1$, since all tested primes exceed $q$. Thus one factor $u(q)$, not $u(q)^r$, converts to a reduced-class average.

\begin{lemma}\label{lem:integer}
For uniform $q$ in the dyadic block, $Z_q/(J u(q)^{r-1})\to1$ in probability.
\end{lemma}
\begin{proof}
Fix $w>2r+2$, $W=\prod_{p\le w}p$ and condition on $q\equiv c\pmod W$, where $c$ need not be reduced. Put $u_w(c)=\prod_{p\le w,\,p\mid c}(1-1/p)^{-1}$. For $d=1,2$ put $e_d=d(r-1)$ and
\[
 K_{w,d}=\prod_{p>w}\left[1+\frac{(1-1/p)^{-e_d}-1}{p}\right]=1+O_r(1/w).
\]
We prove the precise conditional limits, in the order $N\to\infty$ at fixed $w$:
\[
 \E[Z_q^d\mid q\equiv c\pmod W]\ \longrightarrow\ J^d u_w(c)^{e_d}K_{w,d}.\tag{1}
\]

\paragraph{Counting a fixed array of row indices.}
Take $d$ independent copies of the ordered tuples $\mathbf b_\ell$, distinct only within each copy. Write $q=Wn+c$, and use the $dr$ affine prime forms
\[
 a_\ell+b_{\ell i}(Wn+c),\qquad 1\le\ell\le d,\quad 1\le i\le r,
\]
in the $d+1$ variables $(n,a_1,\ldots,a_d)$. In one copy, distinct indices give nonproportional linear parts (coefficient of $a_\ell$ is 1); different copies have independent representative coordinates, even when they use equal row indices. The linear coefficients and constants are $O_{W,\mathcal I}(\log N)$. The body given by $N<Wn+c\le2N$, $0<a_\ell\le Wn+c$ is convex, has volume $\asymp_W N^{d+1}$ and is in a box of side $O(N)$, with all forms positive between $N$ and $O_{\mathcal I}(N\log N)$. Subdivision into a fixed number of equal modulus blocks has the same properties. For any admissible array, Bienvenu \cite[Thm.~1.3]{bien} gives the weighted count as the volume times $\prod_p\beta_p(1+o_{W,r,\mathcal I}(1))$, with uniform \emph{relative} error over the row indices. The local factor is the mean product of the local weights $(1-1/p)^{-1}\1_{p\nmid\text{form}}$ on residue vectors mod $p$. Prime powers can be removed uniformly: for any one form a chosen $n$ has $O((NL)^{1/2}\log(NL))$ exceptional $a_\ell$, the other $d-1$ coordinates have $O(N^{d-1})$ choices, and weights are $O(L^{dr})$. The absolute resulting error is $O(N^{d+1/2}L^{O_r(1)})$ per array. Replacing each prime log by $L$ costs a relative $O_r(\log L/L)$ by the range above.
\par At $p\nmid W$ all arrays are locally admissible: choose $q=0$ and all $a_\ell=1\pmod p$. At a prime $p\mid W$ local inadmissibility means that all possible representatives in at least one copy hit a zero; this makes the actual prime count zero since the tested integers exceed $p$. Thus inadmissible arrays may be included with main product zero. Small absolute errors remain negligible after all $O_{\mathcal I}(L^{dr})$ choices. To insert $q^{-d}$, partition $[N,2N]$ into a fixed number of equal subintervals, use their endpoint bounds for this positive weight, then let the mesh in units of $N$ decrease to zero. Equivalently the weighted real volume is $W^{-1}\int_N^{2N}q^{-d}q^d\,dq=N/W$, which divided by the number of moduli in the cell tends to one. The uniform relative asymptotics on each fixed piece justify these Riemann bounds, since the mean singular product is bounded for fixed $w$, as proved next.

\paragraph{Averaging the local factors over row tuples.}
Given $q\bmod p$, let $\nu_\ell(p)$ be the number of distinct $b_{\ell i}\pmod p$. The conditional local multiplier equals
\[
 \begin{cases}
 (1-1/p)^{-e_d},&q=0\pmod p,\\
 (1-1/p)^{-dr}\prod_{\ell=1}^d(1-\nu_\ell(p)/p),&q\ne0\pmod p .
 \end{cases}\tag{2}
\]
This first averages over the independent $a_\ell\pmod p$. If $p\mid W$ use the appropriate line at $q=c$; if $p\nmid W$ take the mixture with probabilities $1/p$ and $1-1/p$, respectively. For independently uniform $b_{\ell i}\pmod p$, the mean of the second line is one: for each uniform avoided residue, each of the $r$ labels misses it independently. Thus the formal product of the averaged factors is exactly $u_w(c)^{e_d}K_{w,d}$.

Here is the uniform passage to that product. Put $z=(\log L)/4$, so $\prod_{p\le z}p\le L^{1/2}$ eventually. Each integer interval $LI_i$ has $|I_i|L+O(1)$ indices, equidistributed modulo that product with total variation error $O_{\mathcal I}(L^{-1/2})$. Conditioning on no within-copy repeats removes $O_{r,\mathcal I}(1/L)$ probability. The small-factor product is nonnegative and bounded by $(\log(2z))^{O_r(1)}$, so CRT yields its expected independent product up to $o(1)$, for $z>w$. For $p>z$ with no within-copy collision, expression (2) after averaging over $q$ is $1+O_r(p^{-2})$, even if there are cross-copy collisions of labels. Otherwise it is $1+O_r(1/p)$ and is positive for all such $p$. For each nonzero within-copy difference $b-b'$ of size $O_{\mathcal I}(L)$,
\[
 \sum_{p\mid b-b',\,p>z}\frac1p\le\frac{\log(O_{\mathcal I}(L))}{z\log z}=o(1).
\]
There are finitely many such pairs. Hence the entire tail is multiplicatively $1+o(1)$ uniformly over all distinct tuples, and the product of averaged factors has convergent tail as well. The number of row arrays divided by $L^{dr}$ tends to $J^d$, proving (1).

\paragraph{Unconditioning.}
It follows from (1) that
\[
 \limsup_N\E\left[\left(\frac{Z_q}{J u_w(c)^{r-1}}-1\right)^2\biggm|q\equiv c\pmod W\right]
 \ll_r 1/w
\]
with an absolute-in-$c,w$ implied constant for fixed $r$. First average over the finitely many residue cells. Also
\[
 \E\log(u(q)/u_w(q))\ll \sum_{w<p\le2N}\frac1p\left(\frac1p+\frac1N\right)
 \ll 1/w+o(1).
\]
Markov's inequality, followed by $w\to\infty$, gives the lemma. The same estimate from $p=2$ gives $\E\log u(q)=O(1)$, so $u(q)$ is tight in probability (not uniformly bounded for all moduli).
\end{proof}
\subsection{Returning to the first-prime scale}
First attach to row $b\ge1$ the temporary location $u(q)b/L$. For the reduced-class $r$-th factorial measure, evaluated on a fixed rectangle of bounded intervals $I_i$, its mass is $u(q) Z_q$ but with the index intervals replaced by $L I_i/u(q)$ in the definition. On $1\le u(q)\le B$, subdivide the possible $u(q)$ into finitely many bins of mesh $\eta$. In each bin place fixed inner and outer index intervals around every $I_i/u(q)$, by taking the intersection and union over the bin (endpoint choices have no asymptotic effect). For small enough mesh the inner lengths stay positive. The corresponding raw factorial sums squeeze the desired one term by term, even for overlapping rectangles, because within each tuple indices remain distinct. Lemma~\ref{lem:integer} holds in probability jointly for this finite collection. The bounds after multiplication by $u(q)$ tend, up to $O_{B,r,\mathcal I}(\eta)$ on this event, to $\prod_i|I_i|$ since the lengths are $|I_i|/u(q)+O_{\mathcal I}(\eta)$. Let $\eta\to0$ and then $B\to\infty$, using tightness. Hence each fixed factorial rectangle mass in these temporary coordinates converges in probability, over moduli, to its Lebesgue volume.

The physical location of $a+bq$ is $u(q)(b+a/q)/\log q$. On $u(q)\le B$ and any bounded window this differs uniformly from the temporary location by $O_{B,\mathcal I}(1/L)$ since $\log q=L+O(1)$. Squeezing with fixed expanded and contracted intervals transfers the factorial rectangle convergence in probability for rows $b\ge1$ to the true coordinates. Near zero use one-sided intervals; there is no point at zero. For the distribution of the process we may delete row zero: the proportion of reduced classes having such a prime is at most $\pi(q)/\ph(q)\ll u(q)/\log q=o(1)$ even through \emph{all} moduli (use $u(q)\ll\log\log q$). This bounds total variation of the laws with and without this one possible point. We do not require a high-moment bound for the deletion.

Take the countable family of rational interval rectangles of all orders. Convergence in probability on dyadic blocks gives a single density-one set of moduli on which every one of the rectangle masses for the row-zero-deleted process converges: on block $[2^n,2^{n+1})$ enforce the first $j_n$ tests with error $1/j_n$, choosing $j_n\to\infty$ slowly enough that the discarded block fraction tends to zero. The complement has natural density zero by summing dyadic blocks, also for intermediate endpoints. The mixed ordinary moments of counts on disjoint rational intervals converge now by their finite expansions in mixed falling factorials, to those of independent Poisson variables. These laws have finite exponential moments and so are moment-determinate; the method of moments yields joint distributional convergence. First moments on rational compact windows imply tightness of point measures in the vague topology, and the rational half-open intervals determine the unit Poisson process, whose deterministic boundaries have zero intensity. We obtain full process convergence, including arbitrary bounded intervals by approximation. Finally ordinary convergence of ordered arrivals follows from process convergence on compact windows where the limiting points are simple and avoid the boundary; the probability of fewer than $k$ arrivals before $T$ for the limiting process tends to zero as $T\to\infty$. Subtracting consecutive ordered positions gives the independent exponential waiting times. This proves Theorem~\ref{thm:integer}.

\paragraph{Explicit waiting-time formulation.}
Let $p_{\min}(a;q)$ be the least prime in the reduced class of $a$
modulo $q$. On the set $\mathcal Q$ of Theorem~\ref{thm:integer},
for every fixed $t\ge0$,
\[
 \frac{\#\{a\in(\mathbb Z/q\mathbb Z)^\times:
       p_{\min}(a;q)>t\varphi(q)\log q\}}{\varphi(q)}
 \longrightarrow e^{-t}.
\]
This is the Poisson probability of no point in $[0,t]$; the limit
has no atom at either endpoint.

\section{The remaining process at prime moduli with a prescribed pattern}

An integer set $V$ is admissible if $\nu_p(V)=\#(V\bmod p)<p$ for every prime $p$. Put
$\Sg(V)=\prod_p(1-\nu_p(V)/p)(1-1/p)^{-|V|}$, with $\Sg(\emptyset)=1$. This converges positively for any fixed admissible $V$.
\begin{theorem}\label{thm:prime}
Fix $k\ge0$ and $B<\infty$. For each dyadic $N=2^j$, specify one admissible set $V_N\subset[0,B\log N]\cap\mathbb Z$ of cardinality $k$. There is a set of relative density one among prime moduli (possibly depending on this sequence) such that, as primes $q$ tend to infinity in that set, with $N<q\le2N$, the following hold:
\begin{enumerate}
\item The number $A(q)$ of $1\le a<q$ with all $a+vq$ prime for $v\in V_N$ satisfies
$A(q)\sim(q-1)\Sg(V_N)/(\log q)^k$.
\item Choose $a$ uniformly among these $A(q)$ representatives. Place a point at $p/((q-1)\log q)$ for every prime $p\equiv a\pmod q$ \emph{other than} the $k$ prescribed ones. This process converges in distribution to a Poisson process of unit intensity on $[0,\infty)$.
\end{enumerate}
\end{theorem}
Here the whole dyadic tuple is chosen before the modulus is sampled; we do not claim arbitrary pointwise selection of $V_q$ from all log-sized windows. For fixed $V$ the set of good moduli may be chosen simultaneously for all fixed admissible integer tuples with nonnegative entries by a countable diagonal argument. If $v_i/(\log N)\to t_i$ for an ordering of a moving tuple, the full conditional process has limit equal in law to a unit Poisson process plus the deterministic atoms $\sum_i\delta_{t_i}$ (multiplicities allowed).
\par For example take $V=\{0,2\}$. For almost every prime $q$ in relative density, the number of $a$ for which $a,a+2q$ are both prime and \emph{no other} prime in that class lies below $H(q-1)\log q$ is asymptotic to
\[
 e^{-H}\Sg(\{0,2\})\,q/(\log q)^2
\]
for every fixed $H>0$. The prime difference $2q$ here varies: this is not a twin prime theorem. Leung's unpinned every-modulus conditional law \cite{lepp} and Theorem~\ref{thm:integer} do not imply a microscopic, vanishing-probability conditioning law.

As in the classical singular-series averaging method for
Poisson limits \cite{gal}, one sums over the free row indices;
the prescribed rows here remain fixed within each block.

\subsection{Replicated row moments on prime moduli}
Fix bounded intervals $I_1,\ldots,I_r\subset[0,\infty)$ of positive lengths ($r=0$ allowed), $L=\log N$, $J=\prod |I_i|$. Define
\[
 Z_q=\frac1{q-1}\sum_{a=1}^{q-1}\ \sum_{\mathbf b}^{*}
       \prod_{v\in V_N}\1_{\mathbb P}(a+vq)\prod_{i=1}^r\1_{\mathbb P}(a+b_iq),
\]
where $b_i$ ranges over $L I_i\cap\mathbb Z_{\ge0}$ and the star requires all $b_i$ distinct and outside $V_N$. If $r=0$ the inner sum consists of the empty tuple. Let $S=\Sg(V_N)$. Uniformly over the fixed-size admissible base tuple, $S\gg_k1$: for $p>2k$ each factor is at least $(1-k/p)(1-1/p)^{-k}=1+O_k(p^{-2})$, and the finitely many smaller primes have nonzero possibilities only.
\begin{lemma}\label{lem:prime}
For $q$ uniform among primes in $(N,2N]$, $L^k Z_q/(J S)\to1$ in $L^2$, including $r=0$.
\end{lemma}
\begin{proof}
Compute the first and second moments, with $d=1,2$ copies of the $a$ coordinate, at a shared modulus. For each choice of free row tuples (within-copy distinctness; arbitrary overlap across copies), use the following $1+d(k+r)$ affine forms in the $1+d$ variables:
\[
 q,\qquad a_\ell+hq\quad (h\in V_N\cup\{b_{\ell i}:i\le r\},\ 1\le\ell\le d).
\]
All nonconstant linear parts are pairwise nonproportional, including the initial form $q$, because distinct rows at a shared representative differ in their $q$ coefficient and the other representatives are separate variables. Coefficients are $O_{B,I}(\log N)$. The convex body $N<q\le2N$, $N/L^A\le a_\ell<q$ has volume $\asymp N^{1+d}$ and bounded positive values up to $O(NL)$, also after restricting $q$ to any fixed positive-fraction subinterval. Choose $A>4(k+r+1)$. The discarded representatives have negligible contribution even to the normalized replicated moment: for a prime $q$, at most $O_d(L^{-A})$ of the representative vectors are discarded, there are $O(L^{dr})$ row choices, and normalization multiplies by at most $O_k(L^{dk})$. This cutoff makes all prime logs $(1+O_A(\log L/L))L$, including row zero.

Apply Bienvenu \cite[Thm.~1.3]{bien} on this body for each admissible array. It gives a relative $o(1)$ error uniform over these systems, not an absolute error for the possibly large singular series. The prime-only version follows by discarding higher prime powers: for a non-$q$ form there are $O((NL)^{1/2}\log(NL))$ possible representative values at any fixed $q$, and for the form $q$ there are $O(\sqrt N\log N)$ exceptional $q$ choices. Thus the count error even with all log weights is $O(N^{d+1/2}L^{O_{k,r}(1)})$ per tuple. If the local product vanishes, since admissibility fails at some $p\le k+r$ in one copy, then a prime-only solution could use either $q=p$ or some $a_\ell+hq=p$; both are impossible in the truncated body for large $N$. The cases of the sole form $q$ or one anchor and no other rows cause no obstruction. To normalize the representatives by $(q-1)^{-d}$ partition the dyadic interval into finitely many equal bins, bound the weight by its values at the bin ends, and then shrink their relative mesh. Its untruncated weighted real volume is $(1+o(1))N$ because each real fibre has volume $q^d$; the truncation alters this by $O_d(NL^{-A})$. Divide by the number of primes in $(N,2N]$, $(1+o(1))N/L$ by PNT. The one prime log factor $L$ coming from the added $q$ form cancels exactly. These Riemann bounds will apply after normalization by $S^d$ because the mean singular product divided by $S^d$ is bounded, as shown below.

\paragraph{The exact local decoupling.}
At each $p$, the first local weight is $(1-1/p)^{-1}\1_{q\ne0\ (p)}$, of mean one. For a nonzero $q\pmod p$, multiplication by $q$ permutes the residues of each row set. Conditional on this $q$, the independent $a_\ell$ avoid those rows with probabilities $1-\nu_p(V_N\cup\mathbf b_\ell)/p$. None depends on the actual nonzero value of $q$. Consequently the entire singular product for the prime-modulus system is \emph{exactly}
\[
 \prod_{\ell=1}^d \Sg(V_N\cup\mathbf b_\ell).
\]
Independence of row choices across replicas now makes their row-average factorize exactly, unlike at a general integer modulus. We next verify, uniformly for all admissible base tuples permitted in the theorem,
\[
 \frac{1}{\#\{\mathbf b\}^{*}}\sum_{\mathbf b}^{*}
       \frac{\Sg(V_N\cup\mathbf b)}{S}=1+o(1).\tag{1}
\]
The count of free arrays in one copy is $(J+o(1))L^r$, with the usual empty tuple at $r=0$. At $p$, the ratio factor is
\[
 \frac{1-\nu_p(V_N\cup\mathbf b)/p}{1-\nu_p(V_N)/p}(1-1/p)^{-r}.
\]
For independent uniform $b_i\pmod p$ its mean is one: given any residue avoiding the base tuple, the probability the free rows all miss it is $(1-1/p)^r$. Up to $z=(\log L)/4$ all the integer intervals equidistribute mod the primorial $P\le\sqrt L$ with total variation error $O_{I,r}(P/L)$. Removing within-copy duplicates and intersections with the base costs $O_{k,I,r}(1/L)$ in probability. The truncated nonnegative product is at most $(\log(2z))^{O_r(1)}$, so its mean is $1+o(1)$ by CRT. For $p>z$ the local ratio is $1+O_{k,r}(p^{-2})$ unless some free-free or free-pinned difference is zero mod $p$, when it is $1+O_{k,r}(1/p)$. Pinned-pinned collisions alone already cancel in the ratio. The sum of reciprocals $p>z$ dividing any nonzero row difference $n$, $|n|\ll_{B,I} L$, is at most $\log(2+|n|)/(z\log z)=o(1)$. Thus the entire tail is multiplicatively $1+o(1)$ on distinct tuples, proving (1). No pointwise bound above on $S$ was needed.
\par For fixed finite bin mesh the prime-pattern estimate on each bin, followed by (1) and the weight squeeze, now gives
\[
 \E_{q\ {\rm prime}} Z_q^d=(1+o(1))(L^{-k}J S)^d,
 \qquad d=1,2 .
\]
More precisely the relative limsup and liminf differ from 1 by at most $O_d(\text{mesh})$ before decreasing the mesh. Expanding the square proves the lemma.
\end{proof}

\subsection{Point processes and the conditioning event}
With $r=0$, $Z_q=A(q)/(q-1)$ and Lemma~\ref{lem:prime} proves the first assertion in probability on each dyadic block of primes. With $r>0$, dividing by this $r=0$ quantity gives the factorial rectangle masses of the conditioned remaining process in temporary coordinates $b/L$. For any fixed rectangle these approach $J$ in probability. The actual location $(a+bq)/((q-1)\log q)$ differs from $b/L$ by $O_I(1/L)$ on compact windows, uniformly in $a$. Squeezing by fixed slightly smaller and larger rectangles transfers convergence to the true coordinates. All estimates included row zero from the start; it has not been discarded by appealing to an unconditional probability bound under a rare conditioning.
\par On the dyadic blocks, enforce simultaneously the anchor count and a slowly increasing number of rational-box tests in all dimensions with tolerances decreasing to zero, discarding a vanishing fraction of primes in each block. This standard diagonal selection gives one set of relative natural density one (the number of primes in consecutive dyadic blocks is comparable by PNT). On it all these limits hold, and $A(q)>0$ eventually. On disjoint bounded rational intervals, the factorial moment limits are exactly the independent Poisson ones; ordinary joint moments follow by the falling-factorial expansion, and independent Poisson laws are moment determinate. First-moment bounds on compact windows give tightness in the vague topology, and rational intervals determine the limiting point process. This proves the second assertion of Theorem~\ref{thm:prime} and the example by the limiting Poisson void probability. The deterministic anchor locations, if restored, differ from $v_i/L$ by $O_B(1/L)$; addition of their complete finite measures gives the stated full-process limit.

\begin{thebibliography}{9}
\bibitem{bien} P.-Y. Bienvenu, \href{https://www.impan.pl/en/publishing-house/journals-and-series/acta-arithmetica/online/92171/a-higher-dimensional-siegel-walfisz-theorem}{A higher-dimensional Siegel--Walfisz theorem}, Acta Arith. 179 (2017), 79--100; \href{https://arxiv.org/abs/1607.06625}{preprint}, Theorem 1.3 and Example 1.1.
\bibitem{lecounts} S.-K. Leung, \href{https://arxiv.org/abs/2402.07941}{Moments of primes in progressions to a large modulus}, Forum Math. 38 (2026), 555--574.
\bibitem{lepp} S.-K. Leung, \href{https://arxiv.org/abs/2407.05380}{Pseudorandomness of primes at large scales}, Q. J. Math. 76 (2025), 251--263.
\bibitem{gal} P. X. Gallagher, \emph{On the distribution of primes in short intervals}, Mathematika 23 (1976), 4--9 (conditional ordinary-prime Poisson result, singular-series averaging).
\end{thebibliography}
\end{document}
